\section{FHE takeaways}

\begin{enumerate}
\item
  A \emph{fully homomorphic encryption} protocol allows Bob to compute
  some function \(f(x)\) for Alice in a way that Bob does not get to
  know \(x\) or \(f(x)\).
\item
  The hard problem backing known FHE protocols is the \emph{learning
  with errors (LWE)} problem, which comes down to deciding if a system
  of ``approximate equations'' over \({\mathbb{F}}_{q}\) is consistent.
\item
  The main idea of this approach to FHEs is to use ``approximate
  eigenvalues'' as the encrypted computation and an ``approximate
  eigenvector'' as the secret key. Intuitively, adding and multiplying
  two matrices with different approximate eigenvalues for the same
  eigenvector approximately adds and multiplies the eigenvalues,
  respectively.
\item
  To carefully do this, we actually need to control the error blowup
  with the \emph{flatten} operation. This creates a \emph{leveled FHE}
  protocol.
\end{enumerate}
